\documentclass[11pt]{article}
\usepackage[margin=1in]{geometry}
\usepackage{amsmath,amssymb,amsthm}
\usepackage[T1]{fontenc}
\usepackage[colorlinks=true,linkcolor=blue,citecolor=blue,urlcolor=blue]{hyperref}
\usepackage{microtype}
\hypersetup{pdftitle={Focal pedal area ratios of closed elliptic billiards},pdfauthor={Alec Kriebel},pdfsubject={Signed focal pedal area equality and its phase-independent multiplier},pdfkeywords={elliptic billiards, Poncelet polygons, pedal area, Jacobi elliptic functions}}
\newtheorem{theorem}{Theorem}
\newtheorem{lemma}{Lemma}
\newcommand{\sn}{\operatorname{sn}}
\newcommand{\cn}{\operatorname{cn}}
\newcommand{\dn}{\operatorname{dn}}
\newcommand{\Area}{\mathcal A}
\title{Focal pedal area ratios of closed elliptic billiards}
\author{Alec Kriebel\\\small\href{https://orcid.org/0009-0001-9320-500X}{ORCID: 0009-0001-9320-500X}}
\date{October 4, 2026\quad (version 1.0)}
\begin{document}
\maketitle
\begin{abstract}
Reznik, Garcia and Koiller observed equality between the two focal pedal
area ratios of a closed elliptic billiard and its outer tangent polygon.
We prove this equality for every primitive convex or star orbit with a
nondegenerate elliptical caustic. More precisely, the outer and original
signed pedal areas are proportional by a positive constant independent of
phase and common to both foci. The proof identifies a common pair of simple
poles for the two cyclic area traces and matches their residues on a compact
torus. It reconstructs previously used local cancellation arguments within
established elliptic-function methods. We also give an exact triangular
certificate showing that the common focal ratio need not be constant over
the billiard family. Known even and triangular cases and classical methods
are credited explicitly. This is an extensively AI-assisted, unrefereed
research note, with a reproducible verification supplement.
\end{abstract}

\section{The recorded question and its precise resolution}
Let $E$ be a noncircular ellipse with semiaxes $a>b>0$ and foci
$F_\pm=(\pm c,0)$, where $c^2=a^2-b^2$. Fix a strictly nested,
nondegenerate confocal \emph{elliptical} caustic. Let
$P_0,\ldots,P_{N-1}$ be a primitive closed billiard orbit, $N\geq3$,
listed in traversal order. The outer polygon has vertices at successive
boundary-tangent intersections and sidelines equal to those tangents.
For a focus $F$, the pedal polygon consists of its perpendicular feet on
the ordered supporting lines. Every area in this note is the signed sum
\begin{equation}\label{eq:area}
 \Area(X_0,\ldots,X_{N-1})
 =\frac12\sum_{j=0}^{N-1}\det(X_j,X_{j+1}),\qquad X_N=X_0.
\end{equation}
Write $A_\pm$ for the original chord-pedal areas and $B_\pm$ for
the outer tangent-pedal areas.

The displayed observation in \cite{RGK} is
\begin{equation}\label{eq:E}
 \frac{A_+}{A_-}=\frac{B_+}{B_-}.
\end{equation}
It appears as k605 in arXiv:2004.12497v1, Table~5, p.~6;
as k606 in v11, Table~7, p.~9; and as k607 in the journal
Table~7, p.~349. The observation dates to April 2020.
The journal already explains the primitive-even case by central symmetry.
Its invariant-table framing, and the explicit phase-constancy assertion in
the imported problem record (AMR-050-0034), require a distinction:
\eqref{eq:E} is true, but the common ratio in it need not be constant.

\begin{theorem}\label{thm:main}
For the fixed ellipse and caustic there is a real constant $C_0>0$,
independent of the phase of the closed billiard, such that
\begin{equation}\label{eq:M}
 B_+(w)=C_0A_+(w),\qquad B_-(w)=C_0A_-(w).
\end{equation}
For positive traversal all four signed areas are strictly positive.
Thus \eqref{eq:E} holds with nonzero denominators. The statements hold
for every admitted primitive coprime star turning number, as well as
convex orbits. Reversal and positive integral repetitions preserve
the identities and nonvanishing; reversal negates all four signed areas.
The common focal ratio need not be phase constant: Section~5 gives
an exact counterexample.
\end{theorem}

This scope excludes hyperbolic or collapsed caustics, two-period diameters,
unsigned sums of lobe areas and projections clamped to side segments.
The circular case is elementary and separate: a regular star of turning
number $0<\tau<N/2$ has coincident foci, ratios equal to one, and
$C_0=\sec^2(\pi\tau/N)$.

The proof below uses Stachel's canonical coordinates \cite{Stachel} and
classical Jacobi identities \cite{DLMF}. Matching periods and principal
parts followed by Liouville's theorem is established methodology
\cite{KLS,Roitman}. The specialized local cancellations were already used
in this repository's original-pedal and outer-pedal arguments
\cite{Campaign}; they are rederived here, without importing those
arguments' parity-restricted final conclusions. No new general method or
independent discovery of the original observation is claimed.

\section{Coordinates and the actual perpendicular feet}
Uniform scaling normalizes the caustic semiaxes to $1,k'$, where
$0<k<1$ and $k'=\sqrt{1-k^2}$; the foci become $(\pm k,0)$.
All Jacobi functions below have modulus $k$. Put
$K=K(k)$ and $K'=K(k')$. For the positive primitive turning number
$0<\tau<N/2$, with $\gcd(\tau,N)=1$, set
\begin{equation}\label{eq:canonical}
 v=\frac{2K\tau}{N},\quad \delta=2v,\quad
 a=\frac{\dn v}{\cn v},\quad b=\frac{k'}{\cn v},\quad
 P(u)=(-a\sn u,b\cn u).
\end{equation}
Here $0<v<K$, $a>1$, $b>k'$ and $a^2-b^2=k^2$.
Stachel's Theorem~4.3 and equation~(4.9) \cite{Stachel} parametrize
the billiard vertices by $P(w+j\delta)$.
Jacobi addition gives the chord through $P(u-v),P(u+v)$ and the
tangent at $P(u)$, respectively, as
\begin{equation}\label{eq:lines}
 (-\sn u,\cn u/k')\cdot X=1,\qquad
 (-\sn u/a,\cn u/b)\cdot X=1.
\end{equation}
The first line is tangent to the caustic. The midpoint parameter $u$
is essential: the chord numbered $j$ has parameter $w+v+j\delta$.

The projection of $F_+=(k,0)$ onto these lines is
\begin{align}\label{eq:q}
 q(u)&=\left(\frac{k-\sn u}{1-k\sn u},
                  \frac{k'\cn u}{1-k\sn u}\right),\\
 Q(u)&=\left(\frac{a(k-a\sn u)}{a-k\sn u},
                  \frac{ab\cn u}{a-k\sn u}\right).\label{eq:Q}
\end{align}
These formulas follow by substituting each normal $n$ into
$F_++(1-n\cdot F_+)n/(n\cdot n)$, using
$\cn^2u+\sn^2u=1$ and $a^2-b^2=k^2$.
For real $u$ the denominators are at least $1-k>0$ and $a-k>0$.
Consecutive outer tangents intersect finitely: antipodal parameters
have difference $2K$, whereas $0<\delta<2K$.

In the complex continuation the dot product is bilinear. Define
\begin{align}\label{eq:traces}
 A(w)&=\frac12\sum_{j=0}^{N-1}
  \det\bigl(q(w+v+j\delta),q(w+v+(j+1)\delta)\bigr),\\
 B(w)&=\frac12\sum_{j=0}^{N-1}
  \det\bigl(Q(w+j\delta),Q(w+(j+1)\delta)\bigr).\nonumber
\end{align}
These are precisely $A_+(w),B_+(w)$; changing the outer side starting
index only cyclically permutes the summands.
Central inversion is $u\mapsto u+2K$. Projection commutes with it,
and inversion preserves signed area, so
\begin{equation}\label{eq:focus}
 A_-(w)=A(w+2K),\qquad B_-(w)=B(w+2K).
\end{equation}

\section{A common pair of simple poles}
Cyclic reindexing makes $\delta$ a period of both traces;
$4K$ is also a period. Since $\gcd(\tau,N)=1$, integer combinations give
\begin{equation}\label{eq:L}
 L=4K/N
\end{equation}
as a common real period. The imaginary Jacobi shift fixes $\sn$
and negates $\cn$, hence reflects each foot map in the horizontal axis:
\begin{equation}\label{eq:anti}
 A(w+2iK')=-A(w),\qquad B(w+2iK')=-B(w).
\end{equation}
Thus $A$ and $B$ are meromorphic on the compact torus
$\mathbb C/(L\mathbb Z+4iK'\mathbb Z)$.

\begin{lemma}\label{lem:poles}
Modulo this lattice the only possible poles of either trace are
\begin{equation}\label{eq:poles}
 K-v+iK',\qquad K-v+3iK',
\end{equation}
and each has order at most one.
\end{lemma}
\begin{proof}
At a common Jacobi pole, each map \eqref{eq:q}--\eqref{eq:Q}
is a quotient of functions with at most a simple pole and a nonzero
leading denominator coefficient. Both components therefore extend
holomorphically there.

Put $r=K+iK'$. On the $\sn$ torus with periods $4K,2iK'$,
$\sn u=1/k$ has the unique double root $r$. The quarter-shift formulas are
\begin{equation}\label{eq:quarter}
 \sn(r+z)=\frac{\dn z}{k\cn z},\qquad
 \cn(r+z)=-\frac{i k'}{k\cn z}.
\end{equation}
Consequently $q(r+z)$ is even in $z$, and
\begin{equation}\label{eq:germ}
 q(r+z)=\frac{2}{k}(1,i)z^{-2}+O(1).
\end{equation}
At a singular original-foot vertex, its neighbors
$q(r+z\pm\delta)$ are holomorphic: their parameters are not another
root modulo $4K$ since $0<\delta<2K$.
This includes a possible common Jacobi pole at a neighbor, already
shown removable. Its two incident determinants combine into
\begin{equation}\label{eq:originalcancel}
 \tfrac12\det\bigl(q(r+z),q(r+z+\delta)-q(r+z-\delta)\bigr).
\end{equation}
The second vector is odd and holomorphic by the evenness in
\eqref{eq:quarter}, so it is $O(z)$. Equation~\eqref{eq:germ}
then bounds \eqref{eq:originalcancel} by a simple pole.
Primitivity makes all $j\delta$ distinct modulo $4K$; thus exactly
one original-foot vertex can be singular at that phase. All other
area terms are holomorphic.

For the outer map, the nonremovable candidates solve $\sn u=a/k$.
By \eqref{eq:quarter} the roots are $r-v,r+v$; the degree-two
property of $\sn$ leaves no others on its torus. They are distinct
simple roots for $0<v<K$. At both, $\cn u=-ib/k$, and the numerator
vector in \eqref{eq:Q} is
$-ab^2(1,i)/k$. The two denominator derivatives are opposite nonzero
values of magnitude $b^2\sn v$. Thus the residues at $r-v,r+v$
are opposite collinear vectors (indeed $-R,R$ with
$R=a(1,i)/(k\sn v)$).

The roots differ by $\delta$, so two adjacent outer-foot vertices
are singular simultaneously. Their common edge has double-pole
coefficient $\det(-R,R)=0$. Every remaining incident edge has at
most one singular endpoint and therefore at most a simple pole.
There are exactly two such singular vertices: primitivity prevents
additional coincidences, and $N=2$ is excluded. The cyclic closing
edge obeys the same computation, including when $N=3$.

The original phase candidates are $r-v-j\delta$; the outer
candidates are $r-v-j\delta$ and $r+v-j\delta$.
Since $\delta$ is an integer multiple of $L$ and $2v=\delta$,
all reduce to $r-v$ modulo $L$. Their translate by $2iK'$ gives
the second point in \eqref{eq:poles}. No other poles remain.
\end{proof}

\section{Positivity and residue matching}
The positivity needed here is signed positivity, including for stars.
Consider either ellipse with semiaxes $R,S$ equal to $1,k'$ or $a,b$.
Its outward normal in the chosen parametrization is
$n(u)=(-\sn u/R,\cn u/S)$. On the real axis,
\begin{equation}\label{eq:normal}
 \det(n(u),n'(u))=\frac{\dn u}{RS}>0.
\end{equation}
A continuous lift of its direction angle therefore increases strictly,
and its increase over $2K$ is exactly $\pi$. Since
$0<\delta<2K$, consecutive normals have direction increments
strictly between $0$ and $\pi$.

The focus is strictly inside each ellipse. For its tangent with normal
$n$, the foot vector relative to the focus is
$(1-n\cdot F_+)n/(n\cdot n)$, a positive multiple of $n$.
Every consecutive determinant of these vectors is positive. The same
holds for the closing edge using the lifted increment;
$N\delta=4K\tau$ closes the directions after $\tau$ full turns.
Translating the origin to the focus leaves \eqref{eq:area} unchanged.
It follows that $A(w)>0$ and $B(w)>0$ at every real phase.
Equation~\eqref{eq:focus} gives positivity for the other focus too.

A nonzero trace with the pole bound of Lemma~\ref{lem:poles}
must have genuine simple poles at both listed points. If one were
absent, \eqref{eq:anti} would remove the other. A holomorphic function
on the compact torus is constant, and \eqref{eq:anti} would force
that constant to be zero, contrary to positivity.

At the first pole choose
$C_0=\operatorname{Res}(B)/\operatorname{Res}(A)$; the denominator
residue has just been proved nonzero. The difference $B-C_0A$
has no pole there. Its antiperiodicity removes the pole at the second
point as well. Lemma~\ref{lem:poles} makes it holomorphic everywhere,
hence constant, hence zero by \eqref{eq:anti}.
On the real axis $C_0=B(w)/A(w)>0$; it is real and independent of phase.
Applying this identity at $w+2K$ proves the identical multiplier at
$F_-$, then permits division to obtain \eqref{eq:E}.
Reversal negates each area and a repeated traversal multiplies each
by its positive repetition count. This proves the positive assertions
of Theorem~\ref{thm:main}.

\section{An exact nonconstant focal ratio}
Consider $x^2/21+y^2/16=1$ with foci $(\pm\sqrt5,0)$ and caustic
squared semiaxes $189/25,64/25$. A primitive positively oriented
triangular billiard is
\begin{equation}\label{eq:triangle}
 (\sqrt{21},0),\quad(-3\sqrt{21}/5,16/5),\quad
 (-3\sqrt{21}/5,-16/5).
\end{equation}
Its lines are $x=-3\sqrt{21}/5$ and
$2x/\sqrt{21}\pm y=2$. The caustic support test gives
$(189/25)(2/\sqrt{21})^2+64/25=4$ on each oblique side;
the vertical side is also tangent. The contact points lie on the
segments and direct normalized-vector reflection verifies the billiard
law. These exact checks are included in the supplement.
Perpendicular projection and \eqref{eq:area} give
\begin{equation}\label{eq:exactareas}
 A_\pm=\frac{84(7\sqrt{21}\pm\sqrt5)}{625},\qquad
 B_\pm=\frac{7(7\sqrt{21}\pm\sqrt5)}{10}.
\end{equation}
The common multiplier is $125/24$, whereas the common focal ratio is
$(7\sqrt{21}+\sqrt5)/(7\sqrt{21}-\sqrt5)>1$.
Central inversion gives a different phase of the same fixed triangular
family, preserves orientation, and exchanges the foci. Its ratio is
therefore the reciprocal, strictly less than one. This proves the
negative assertion of Theorem~\ref{thm:main} without a numerical test.
It corrects the constancy interpretation of the source and imported
record; it does not refute the displayed equality.

\section{Prior cases, verification and limitations}
The negative result is also a short consequence of classical inputs.
For a signed triangle $T$ with circumcenter $O$ and radius $R$,
the Querret--Sturm pedal formula is
\begin{equation}\label{eq:querret}
 D_T(F)=\frac{[T]}{4R^2}\bigl(R^2-|F-O|^2\bigr) \cite{Querret}.
\end{equation}
Fierobe's Lemma~4.1 \cite{Fierobe} distinguishes the two axial
triangular circumcenters. Applying \eqref{eq:querret} at the foci
and then inverting the triangle gives unequal reciprocal ratios.
Thus no new negative theorem or first correction is claimed.

The triangular positive case is likewise an explicit consequence of
older geometric statements \cite{Book,Center}. The supplement gives
the complete conditional deduction, including the simultaneous
center-coordinate relation that a locus shape alone would not imply.
For $X=a^2$, $Y=b^2$, $D=\sqrt{X^2-XY+Y^2}$ and
$t=(D-Y)/(X-D)>1$, its multiplier is
$C_0=(t+1)^3/(2t)$. This is credited as an old-results corollary,
not a separately novel triangular result. The main proof does not
rely on those triangular center theorems or on their illustrative
printed examples.

Recent restricted pedal results \cite{Recent} use related trace
methods and different quantities or least-period classes. The
verification supplement records their exact comparison and six
unavailable journal editions. A bounded primary-source audit found
no previous full-domain proof of \eqref{eq:E} or of the stronger
proportionality statement \eqref{eq:M}; finite literature coverage cannot certify global
firstness. In particular, no new even-period multiplier priority or
independent discovery of $C_0$ is asserted.

The supplement contains original portable exact controls, finite
high-precision real and complex diagnostics, and an independent
reflection construction. Computation supplements the proof rather
than establishes it by sampling; the numerical results are not
interval-certified bounds. AI tools were used extensively for proof
development, literature comparison, computation, writing and
adversarial verification. Internal independent AI reviews are not
human peer review. This note has not been refereed, independently
human peer reviewed, or certified by a formal proof assistant.

\begin{thebibliography}{12}
\bibitem{RGK} D. S. Reznik, R. Garcia and J. Koiller,
\emph{Fifty New Invariants of N-Periodics in the Elliptic Billiard},
Arnold Math. J. \textbf{7} (2021), 341--355,
\href{https://doi.org/10.1007/s40598-021-00174-y}{doi:10.1007/s40598-021-00174-y}.
Earlier versions: \href{https://arxiv.org/abs/2004.12497v1}{arXiv:2004.12497v1}
and \href{https://arxiv.org/abs/2004.12497v11}{v11}.
\bibitem{Stachel} H. Stachel,
\emph{On the motion of billiards in ellipses},
European J. Math. \textbf{8} (2022), 1602--1622,
\href{https://doi.org/10.1007/s40879-021-00524-2}{doi:10.1007/s40879-021-00524-2},
Theorem~4.3 and equation~(4.9), p.~1614.
\bibitem{DLMF} NIST Digital Library of Mathematical Functions,
\href{https://dlmf.nist.gov/22.4}{\S22.4} and
\href{https://dlmf.nist.gov/22.8}{\S22.8}, Jacobi periods and addition identities.
\bibitem{KLS} A. Khare, A. Lakshminarayan and U. Sukhatme,
\emph{Local Identities Involving Jacobi Elliptic Functions},
\href{https://doi.org/10.1007/BF02704435}{doi:10.1007/BF02704435} (2004),
\S4.1, equations~(39)--(42);
\href{https://arxiv.org/abs/math-ph/0306028}{arXiv:math-ph/0306028}.
\bibitem{Roitman} P. Roitman, R. Garcia and D. Reznik,
\emph{New Invariants of Poncelet--Jacobi Bicentric Polygons},
Arnold Math. J. \textbf{7} (2021), 619--637,
\href{https://doi.org/10.1007/s40598-021-00188-6}{doi:10.1007/s40598-021-00188-6}.
\bibitem{Querret} Querret,
\emph{D\'emonstration du th\'eor\`eme de g\'eom\'etrie \'el\'ementaire
\'enonc\'e \`a la page 28 du pr\'esent volume},
Annales de Math\'ematiques pures et appliqu\'ees \textbf{14}
(1823--1824), 280--285,
\href{https://www.numdam.org/item/AMPA_1823-1824__14__280_0/}{Numdam primary article}.
\bibitem{Fierobe} C. Fierobe,
\emph{On the circumcenters of triangular orbits in elliptic billiard},
\href{https://arxiv.org/abs/1807.11903v5}{arXiv:1807.11903v5} (22 July 2019), Lemma~4.1.
\bibitem{Book} R. A. Garcia and D. S. Reznik,
\emph{Discovering Poncelet Invariants in the Plane}, IMPA (2021),
ISBN 978-65-89124-43-6, Theorems~2.1 and~2.3.
\bibitem{Center} M. Helman, D. Laurain, R. Garcia and D. Reznik,
\emph{Invariant Center Power and Elliptic Loci of Poncelet Triangles},
\href{https://arxiv.org/abs/2102.09438v4}{arXiv:2102.09438v4} (16 April 2021), \S\S3.4--3.5.
\bibitem{Campaign} A. Kriebel, Math research repository,
original-foot local cancellation,
\href{https://github.com/AlecKriebel/Math/blob/7b7fcd63d5ea83df3a4c29853c4ba34b93d16f07/unsolved_math_prioritization/attempts/5100021/PROOF.md}{commit 7b7fcd6}, \S\S5,7;
outer-foot cancellation,
\href{https://github.com/AlecKriebel/Math/blob/2180d64b648b81ac660a792f7823d0b715f9cf12/unsolved_math_prioritization/attempts/5100033/PROOF.md}{commit 2180d64}, \S\S2--4 (2026).
\bibitem{Recent} A. Ferudun,
\emph{Focal Pedal-Antipedal Area Products in Elliptic Billiards}, v1.0,
\href{https://doi.org/10.5281/zenodo.23079799}{doi:10.5281/zenodo.23079799};
\emph{Steiner-Centroid Pedal Area Ratios in Elliptic Billiards}, v1.0,
\href{https://doi.org/10.5281/zenodo.23089778}{doi:10.5281/zenodo.23089778};
\emph{Outer-Pedal Area Products in Elliptic Billiards}, v1.1,
\href{https://doi.org/10.5281/zenodo.23088516}{doi:10.5281/zenodo.23088516}
(1 October 2026).
\end{thebibliography}
\end{document}
